\begin{statement}{A.31}{Proved}{The stable \(S_3\) fixed subgroups at the prime two}
Let
\[
 C:\ v^2=w^7-w,\qquad
 D:\ z^2=(t^2+4)(t^3+3t),
\]
with the two indicated cusps removed from \(C\), and the unique
cusp removed from \(D\). For every original packet and every
permitted open subgroup \(\Delta_0\) on which the selected lift
is defined,
\[
 P_{C,2}^{\,s_b(\Delta_0)}=I_{C,2},
 \qquad
 P_{D,2}^{\,s_b(\Delta_0)}=I_{D,2}.                 \tag{A.31.1}
\]
The second assertion applies to each reflection quotient.
Together with A.27 and A.30, these statements give the full
fixed-subgroup equality for the stable \(S_3\) genus-three cover
and all its genus-two reflection quotients at every prime.

The replacement for prime-to-two tame monodromy is Fontaine's
monodromy on the semistable realization. It is used on a
crystalline realization of a global orbit span, never as an
operator on an individual etale vector.
\end{statement}
\begin{proof}
\paragraph{The semistable comparison used at a cusp.}
We first specify the comparison input, including its tangential
form. Let \(F/\mathbb Q_p\) be finite. Suppose that a marked
curve has a regular semistable model over \(\mathcal O_F\),
with disjoint marked sections in its smooth locus. Give the
model the log structure of the special fibre and the marked
sections. At an integral nonzero tangent to a marked section,
every finite-dimensional quotient of the rational unipotent
fundamental Lie algebra is semistable. Its \(D_{\rm st}\)
realization is the corresponding logarithmic crystalline
Lie algebra, compatibly with brackets, the lower-central
filtration, and the cusp map \(\mathbb Q_p(1)\to\mathfrak u\).
Its monodromy is the logarithmic Gauss--Manin residue.

Here is the passage from the ordinary to the tangential fibre
in this assertion. The ramified semistable comparison is that
of \cite[Corollary 5.1 and Theorem 5.2]{fghAIK}: it is an
association of the universal unipotent local systems with
logarithmic crystalline objects on the whole marked model,
before taking the fibre. The open-curve scope is explicit in
\cite[Remark 1.10 and Section 2]{fghAIK}.
Here is the literal local period map in this pullback.
Write \(W=W(k)\), choose a uniformizer \(\varpi\) of \(F\),
and set
\[
 \mathcal O=W[[Z]],\qquad Z\longmapsto\varpi,\qquad
 \mathcal O_{\rm cris}
   =\widehat{\operatorname{DP}}_{\,(p,P_\varpi(Z))}
                                      W[[Z]],
\]
as in Section 2.1 of \cite{fghAIK}.
On the lifted smooth chart use \(W[[Z,z]]\), with
horizontal marked divisor \(z=0\), and choose the
tangent with normal coordinate \(1\).
Let \(\overline R\) be the integral Kummer tower of
this chart used in the definition of its period sheaf.
A tangential specialization of this tower gives a
map \(\overline R\to\mathcal O_{\mathbb C_p}\).
It sends the underlying elements \(z^{1/p^n}\) to
zero, but sends their logarithmic generators to
compatible formal generators \(e_h/p^n\), whose
structure map is zero. It consequently induces
both a map of tilt rings and a map of their log
monoids. This is the specialization construction
in \cite[Section 9.25, especially (9.25.4)--(9.25.7)]
{fghOlssonTangential}.

Combine its Witt-vector map with \(z\mapsto0\)
and \(Z\mapsto Z\). By the universal property of
the log divided-power envelope, the result is
a morphism from the AIK local log-period
envelope to
\[
       \bigl(A_{\log}(F),\,
                         M_{\rm vertical}\oplus\mathbb N e_h\bigr).
                                                        \tag{A.31.12}
\]
The horizontal copy from the lifted chart and the
horizontal copy from the tilt both map to \(e_h\).
Thus their \emph{log exactification unit} \(u_h\)
maps to \(1\), and every positive divided power
of \(u_h-1\) maps to zero. This is a map of log
envelopes; no numerical quotient \(0/0\) is taken.
In contrast, the vertical exactification unit
\([\varpi^\flat]/Z\) is retained. The underlying
target ring is exactly the vertical \(A_{\log}(F)\)
of AIK, with its additional horizontal log
generator of structure value zero. The
Eisenstein relation \(P_\varpi(Z)\) and its
divided powers remain in the vertical base.
This explains explicitly why ramification of
\(F\) causes no new horizontal relation.

Frobenius sends both log generators to their
\(p\)-fold multiples and sends every
exactification unit to its \(p\)-th power.
It therefore commutes with (A.31.12).
The tangential splitting makes the formal
\(e_h/p^n\) fixed by the arithmetic \(G_F\)
section, so the map is \(G_F\)-equivariant.
It also commutes with the vertical logarithmic
connection, since \(u_h\) has vertical derivative
zero and the vertical unit is unchanged.
The divided-power filtration is respected.
Inverting the Tate period gives the same map
with target \(B_{\log}(F)\).
Pullback of a crystal along this morphism of
logarithmic enlargements is its evaluation
on the target; hence the AIK coefficient
association pulls back to an association on
the horizontal log point. This is the
base-change identity (9.26.2)--(9.26.4) of
\cite{fghOlssonTangential}, with the vertical
log base retained.

Relative to that vertical base, a unipotent
object at this log point is a finite module
with a nilpotent horizontal residue \(R\).
Nilpotence follows from its geometric
unipotent filtration: \(R\) is zero on each
constant graded quotient. The tangent
retraction evaluates this object on the
chosen \(e_h\), giving its usual tangential
fibre. Equivalently, with
\(\ell_h=\log u_h\), the logarithmic Taylor
identification is the finite exponential
\(\exp(-\ell_hR)\), followed by
\(\ell_h=0\). Integrability says that \(R\)
is horizontal for the vertical connection;
thus this identification commutes with
vertical monodromy. Frobenius compatibility
follows from \(\varphi(\ell_h)=p\ell_h\)
and the corresponding residue relation
obtained from \(d\log(\varphi z)=p\,d\log z\).
On the etale side horizontal inertia is
\(\mathbb Q_p(1)\); its comparison is exactly
the Tate-period calculation of
\cite[Proposition 9.5, Corollary 9.16, and
Section 9.24]{fghOlssonTangential}.
This verifies the arithmetic fibre
association and all needed compatibilities.

The intermediate coefficient ring here is
\(B_{\log}\), not a quotient identified
directly with \(B_{\rm st}\).
The fibre is \(B_{\log}\)-admissible.
\cite[Proposition 2.1]{fghAIK}, namely
Breuil's equivalence, converts this into
\(B_{\rm st}\)-admissibility and identifies
the monodromy with the vertical residue.
This construction uses Olsson's local
log-point calculation and AIK's ramified
coefficient comparison. It does not apply
Olsson's global unramified good-reduction
theorem to a ramified semistable model.

One can take the AIK universal objects based at any ordinary
point through the smooth locus, evaluate them at this log
point, and reconstruct its unipotent group by their tensor
matrix coefficients. These finite Lie quotients are subquotients
of finite sums and tensor products of the associated fibres.
Semistable representations are closed under these operations,
so all such Lie quotients are semistable. The same construction gives the
path torsors and the compatibility with the cusp.
Finally \cite[Lemma 6.2(ii) and the deformation argument in
Section 6.2]{fghAIK} identifies the vertical residue with
Fontaine's \(N\). Deformation to a complex marked semistable
curve identifies it with the logarithm of the usual edge
twists. If a node has equation \(xy=\varpi^e\), its
coefficient is \(e\), up to the single common nonzero
normalization of \(N\). Rational subdivisions add these
coefficients along an arm. This statement is about the
residue operator, and hence gives the edge-twist formulas
in every finite unipotent quotient.

\paragraph{Every cusp packet remains semistable.}
The preceding comparison applies first to an integral
algebraic tangent. It also applies to every packet used
here, after rationalization. Locally the difference from
such a tangent is
\[
                  \kappa\in H^1(G_F,\mathbb Q_p(1)).
\]
Every such class is semistable. Indeed Kummer theory identifies
this finite-dimensional vector space with the rationalized
pro-\(p\) completion of \(F^\times\). A uniformizer and a
basis of principal units generate it over \(\mathbb Q_p\).
Their Kummer extensions split over \(B_{\rm st}\):
the logarithm of the compatible unit roots belongs to
\(B_{\rm cris}\), and the uniformizer logarithm belongs to
\(B_{\rm st}\). The explicit construction is also given in
\cite[Sections 11.7--11.10]{fghOlssonSemistable}; this
Kummer-period calculation works over every finite \(F\).
Thus there is
\[
 b_\kappa\in B_{\rm st}\otimes_{\mathbb Q_p}\mathbb Q_p(1),
 \qquad
                 \kappa(g)=g(b_\kappa)-b_\kappa .
                                                        \tag{A.31.2}
\]
The sign can of course be reversed with the convention for
the Kummer cocycle.

Let \(c:\mathbb Q_p(1)\to\mathfrak u\) be the actual cusp
map. The twisted action is
\(\rho_\kappa(g)=\exp(\operatorname{ad}c(\kappa(g)))\rho(g)\).
Over \(B_{\rm st}\), conjugation by
\(\exp(\operatorname{ad}c(b_\kappa))\), with the sign matching
(A.31.2), identifies this with the untwisted action.
The exponentials are finite in each nilpotent quotient.
This proves semistability for every packet, not just for
packets represented by an algebraic element of \(F^\times\).
It also shows directly why a packet can add a cusp-inner
term to \(N\). We shall allow that term rather than choose
a special packet.

The selected lift to either cover is itself cuspidal after
finite descent. Its ramification index over the original
cusp is three; at \(p=2\) division by three is integral.
Consequently the preceding local argument applies to all
the induced original packets in the statement.

\paragraph{A semistable representation with abelian image is crystalline.}
We give the linear-algebra proof needed below. Let \(W\)
be semistable and suppose that its Galois image is abelian.
The finite-dimensional commutative algebra generated by
that image splits after a finite extension of coefficient
fields. Decompose \(W\) into its generalized-character
blocks. Each character \(\eta\) is a semistable rank-one
subquotient, and hence crystalline: a nilpotent endomorphism
of its rank-one \(D_{\rm st}\) module is zero.
Twist its block by \(\eta^{-1}\). All composition factors
are now the trivial representation. Exactness of
\(D_{\rm st}\) implies that \(\varphi^f\) on this block
is unipotent, where \(f\) is the residue degree of \(F\).
The relation
\[
                       N\varphi^f=p^f\varphi^fN
                                                        \tag{A.31.3}
\]
then forces \(N=0\): it would carry the generalized
eigenspace with eigenvalue \(1\) into the generalized
eigenspace with eigenvalue \(p^{-f}\), which is zero.
Twisting back by the crystalline character preserves
this conclusion. It descends through the coefficient
extension, proving the assertion. This also covers
unipotent abelian representations; no semisimplicity
of the original image is assumed.

\paragraph{A fixed element supplies crystalline leading classes.}
Enlarge a finite number field \(L\) so that the cover,
packet lift, cusps, CM isogenies and marked semistable
models at a place \(v\mid2\) are defined, and
\[
       \Gamma=\overline{[G_L,G_L]}\subset\Delta_0 .
\]
This is possible because \(\Delta_0\) is open in \(\Delta\).
Put \(F=L_v\). If \(u\) is an actual fixed pro-\(2\)
element, its logarithm in a finite rational nilpotent
quotient \(\mathfrak u_n\) is fixed by \(\Gamma\).
The span
\[
 M_n=\left\langle \rho(g)\log(u)_n:g\in G_L
                                            \right\rangle_{\mathbb Q_2}
                                                        \tag{A.31.4}
\]
is also fixed by \(\Gamma\), since \(\Gamma\) is normal
in \(G_L\). Its Galois image is therefore abelian.
The first two paragraphs make \(\mathfrak u_n\), and
hence \(M_n\), semistable locally at \(v\).
The preceding paragraph gives
\[
                         N D_{\rm st}(M_n)=0.           \tag{A.31.5}
\]

Both proper Jacobians have good reduction after this
extension. Thus their \(H_1\) representations are
crystalline. For \(C\), ordinary open homology splits
rationally into the boundary line \(\mathbb Q_2(1)\)
and proper homology: the hyperelliptic involution acts
by \(+1\) on the former and by \(-1\) on the latter.
The projectors \((1\pm h)/2\) are legitimate over
\(\mathbb Q_2\). It follows that all ordinary graded
pieces \(\operatorname{gr}_d\mathfrak u\), free-Lie
constructions from \(H_1\), are crystalline.
After the chosen finite extension the global action
on \(H_1\) is abelian, given by the two CM characters
and, for \(C\), the cyclotomic boundary character.

Suppose that \(u\) has first nonzero degree \(d\), and
let \(W_d\) be the global \(G_L\)-span of its leading
class. For \(n\) sufficiently large, \(M_n\) is in
lower-central degree at least \(d\) and maps onto
\(W_d\). If \(N_r\) is the first nonzero graded
derivation of \(N\), exactness of \(D_{\rm st}\)
and (A.31.5) give
\[
                 D_{\rm cris}(W_d)\subset\ker N_r .
                                                        \tag{A.31.6}
\]
Indeed lift a vector of \(D_{\rm cris}(W_d)\) to
\(D_{\rm st}(M_n)\) and take the component of \(N\)
in degree \(d+r\). All smaller degree components
vanish, so that component is \(N_r\) of the leading
vector.

There is also a genuine CM-stability statement on
the left side of (A.31.6). Every global Galois operator
on \(H_1\) commutes with the local \(G_F\)-action,
because the global homology action is abelian.
It is therefore a morphism of crystalline
representations. The same is true on each free-Lie
graded piece and on \(W_d\). Applying \(D_{\rm cris}\)
shows that \(D_{\rm cris}(W_d)\) is stable under the
induced global CM operators. In particular it lies
in the intersection of \(\ker N_r\) with every one
of its CM translates.
All component maps, polarizations and CM
endomorphisms used in A.27--A.28 are compatible
with crystalline realization. Their displayed
linear formulas consequently hold here after a
coefficient extension splitting the CM field.
This proves the required transport statement
entirely inside the crystalline realization.
It does not identify \(N\) with an operator on
the original \(\mathbb Q_2\)-vector space.

\paragraph{The genus-two curve.}
The pointed stable model is the two-tail model of
A.27, with the cusp on its rational bridge.
Its ordinary-based residue has no degree-one term
and has degree-two term
\[
                           m_1T_1+m_2T_2,
                     \qquad m_1m_2\ne0.                \tag{A.31.7}
\]
Compare it with the actual packet's \(N\).
They have the same affine outer residue. Over a
characteristic-zero coefficient field a unipotent
path torsor has a point, so the two based residues
differ by an inner derivation. Homology is
crystalline, so both have degree zero equal to
zero. Their possible degree-one difference is
\(\operatorname{ad}(w_1)\).
The actual cusp map is crystalline, and thus
\(N(c)=0\). Its leading class is
\(\Omega=\omega_1+\omega_2\), so
\([w_1,\Omega]=0\). The free-Lie centralizer of
\(\Omega\) is its own line, in degree two;
therefore \(w_1=0\).
The degree-two difference is then
\(\operatorname{ad}(w_2)\), with
\([w_2,\Omega]=0\), and hence
\[
                  N_2=m_1T_1+m_2T_2
                                +\lambda\operatorname{ad}(\Omega)
                                                        \tag{A.31.8}
\]
for some scalar \(\lambda\). This allows every
packet-inner ambiguity.

The unbounded kernel calculation (A.27.3) applies
over this characteristic-zero field. If
\(\lambda\ne0\), its kernel is the line \(\Omega\).
If \(\lambda=0\), its kernel is
\(\mathbb L(\omega_1,\omega_2)\).
Choose a global ordinary CM Frobenius away from
two with eigenvalue ratio of infinite order.
The crystalline versions of (A.27.11)--(A.27.15)
give the intersection of this kernel with its
translate as the line \(\Omega\).
Equation (A.31.6) and CM stability therefore imply
that every nonzero leading fixed class is a
multiple of \(\Omega\) in degree two.
Faithfulness of crystalline comparison descends
this containment to etale graded homology.

\paragraph{The genus-three curve.}
Use the three-tail pointed model of A.28.
Its ordinary-based residue has zero first degree
and degree-two part
\[
 N_2(t)=0,\qquad
 N_2|_{V_i}=m_i\operatorname{ad}(\omega_i),
                         \quad m_i\ne0.                \tag{A.31.9}
\]
The same affine outer-residue comparison applies.
Because the actual selected cusp \(c\) is crystalline
and has leading degree-one class \(t\), the possible
inner degree-one term must be
\(\lambda\operatorname{ad}(t)\).
If \(\lambda\ne0\), its kernel is the line \(t\),
and (A.31.6) already gives the desired conclusion.
If \(\lambda=0\), the degree-two inner correction
is \(\operatorname{ad}(w_2)\) with \([w_2,t]=0\);
the centralizer of \(t\) has no degree-two part,
so \(w_2=0\). Thus (A.31.9) is the actual \(N_2\).

The first-core proof of A.28 gives its kernel
\(\mathbb L(t,\omega_1,\omega_2,\omega_3)\).
The component maps have degree one, so the
crystalline CM matrices are exactly (A.28.5).
Their diagonal-plane intersection and its
tensor-algebra consequence give
\[
         D_{\rm cris}(W_d)\subset
              \mathbb L\bigl(t,\Omega\bigr)_d,
                \qquad \Omega=\omega_1+\omega_2+\omega_3 .
                                                        \tag{A.31.10}
\]
The right side is itself the crystalline realization
of the etale subspace generated by the boundary
line and the proper polarization tensor. Taking
the crystalline image of the map to its quotient,
and using faithfulness, descends (A.31.10) to
\(W_d\subset\mathbb L(t,\Omega)_d\).

The first-Johnson nonvanishing established inside
the proof of A.30 is valid over \(\mathbb Q_2\)
as well. Its universal regulator comparison is
rational; the factors \(2\) and \(4\) are invertible
in this field. The explicit primitive functional
is still \(q_*-2(q_0)_*\), and its value is
\[
                      -4\,\kappa_{\,2+\sqrt3}.
                                                        \tag{A.31.11}
\]
The unit \(2+\sqrt3\) has infinite order, so its
rational \(2\)-adic Kummer class is nonzero after
every finite number-field extension. The restriction
to \(G_L\cap\Delta\) is injective on this class:
the abelian quotient has a nontrivial cyclotomic
scalar, so Sah's lemma kills its rational
\(H^1(-,\mathbb Q_2(1))\).
The restricted class stays nonzero on every open
subgroup, since it is then an additive map to a
characteristic-zero vector space.
Packet changes add only a scalar first-Johnson
matrix, killed by this primitive functional.
Hence the actual first Johnson matrix
\(A_\delta\) is nonzero for some
\(\delta\in\Delta_0\), for every packet.

The algebraic implication (A.29.4) uses only
characteristic zero and \(A_\delta\ne0\);
it has no odd-prime hypothesis. Combining it
with (A.31.10) puts every nonzero leading fixed
class on the line \(t\).
This use of A.29 does not invoke its separate
odd-prime integral uniqueness assertion.

\paragraph{Return to the actual pro-\(2\) groups.}
The lower-central filtration of a finitely generated
free pro-\(2\) group is separated, and its graded
pieces are the torsion-free integral free Lie
pieces. On \(D\) the selected cusp has primitive
leading class \(\Omega\) in degree two.
On \(C\) it has primitive class \(t\) in degree one.
Thus the leading coefficient of an actual fixed
element on the indicated line belongs to
\(\mathbb Z_2\). Subtract the corresponding
\emph{actual} cusp power. The remainder is fixed
and has strictly larger first degree; the
graded restrictions above force it to be the
identity. All cusp powers are fixed on
\(\Delta_0\), since its cyclotomic character is
trivial. This proves both equalities (A.31.1)
and completes the assertion for every reflection
quotient.
\end{proof}
